% ═══════════════════════════════════════════════════════════════════════
\section{Déterminer une transformation à partir de points}
% ═══════════════════════════════════════════════════════════════════════

\begin{frame}{La question inverse}
  \small Jusqu'ici : « voici une transformation, que fait-elle à cette
  figure ? ». Maintenant, la question qui se pose \textbf{en pratique} :
  « voici une figure et son image, quelle transformation les relie ? »
  --- le problème que résout un logiciel de retouche photo, ou de
  reconnaissance de formes.
\end{frame}

\begin{frame}{Théorème 1.15 --- Trois points suffisent}
  \begin{Enonce}{Théorème 1.15}
    \small $P_1,P_2,P_3$ \textbf{non alignés}, et $Q_1,Q_2,Q_3$
    quelconques. Il existe une \textbf{unique} application affine $w$
    telle que $w(P_i)=Q_i$, $i=1,2,3$.
  \end{Enonce}
  \small Surnommée \defi{méthode du bulldozer} : on écrit toutes les
  conditions et on pousse tout dans un système linéaire.
  \begin{AuTableau}
    \small Idée : le système à 6 inconnues $(a,b,c,d,e,f)$ se
    \textbf{découple} en deux systèmes $3\times 3$ de même matrice
    \[
      M = \begin{pmatrix} x_1&y_1&1\\x_2&y_2&1\\x_3&y_3&1\end{pmatrix},
      \qquad \dete(M) = 2\cdot\text{aire orientée}(P_1P_2P_3) \neq 0.
    \]
  \end{AuTableau}
\end{frame}

\begin{frame}{Remarque 1.11 --- Une écriture plus compacte}
  \begin{Rem}{Remarque 1.11}
    \small En coordonnées homogènes, tout tient sur une ligne. Avec
    $\widetilde P,\widetilde Q$ les matrices $3\times3$ des points
    (en colonnes, ligne du bas = 1) :
    \[
      \widetilde W = \widetilde Q\,\widetilde P^{-1}.
    \]
    C'est la formule qu'implémentent les bibliothèques d'infographie.
  \end{Rem}
\end{frame}

\begin{frame}{Exemple 1.10 --- Le problème des deux triangles}
  \small $A=(-11,3)$, $B=(1,12)$, $C=(-1,14)$
  $\longmapsto$ $P=(90,-68)$, $Q=(69,-53)$, $R=(83,-59)$.
  \[
    w\col{x}{y} = \mat{-4}{3}{2}{-1}\col{x}{y} + \col{37}{-43}.
  \]
  $\dete(A)=-2$ : renverse l'orientation, double les aires. Ni une
  isométrie, ni une similitude ($A^\top A$ n'est pas multiple de $I$).
  \FigBox[0.36\textheight]{%
    \begin{tikzpicture}[scale=0.2]
      \draw[quadrillage] (-12,-1) grid (3,16);
      \draw[axe] (-12,0) -- (3.6,0) node[etiq, below] {$x$};
      \draw[axe] (0,-1) -- (0,16.6) node[etiq, left] {$y$};
      \fill[fig] (-11,3) -- (1,12) -- (-1,14) -- cycle;
      \node[pt, label={[etiq]left:$A$}] at (-11,3) {};
      \node[pt, label={[etiq]right:$B$}] at (1,12) {};
      \node[pt, label={[etiq]above:$C$}] at (-1,14) {};
    \end{tikzpicture}%
    \hspace{1cm}
    \begin{tikzpicture}[scale=0.2]
      \begin{scope}[shift={(-77,59)}]
        \draw[quadrillage] (68,-69) grid (91,-52);
        \draw[axe] (68,-69) -- (91.6,-69);
        \draw[axe] (68,-69) -- (68,-51.4);
        \fill[figim] (90,-68) -- (69,-53) -- (83,-59) -- cycle;
        \node[ptim, label={[etiq]right:$P$}] at (90,-68) {};
        \node[ptim, label={[etiq]left:$Q$}] at (69,-53) {};
        \node[ptim, label={[etiq]above right:$R$}] at (83,-59) {};
      \end{scope}
    \end{tikzpicture}}
  \begin{center}\small Figure 1.11 --- Le problème du bulldozer.\end{center}
\end{frame}

\begin{frame}{Remarque 1.12 --- Combien de points pour quel groupe ?}
  \begin{Rem}{Remarque 1.12}
    \small Analogues du théorème 1.15 pour les sous-groupes :
    \begin{itemize}
      \item une \textbf{similitude} est déterminée par l'image de
        \textbf{deux} points distincts (à ceci près : deux solutions,
        directe/indirecte) ;
      \item une \textbf{isométrie}, pareil ;
      \item une \textbf{translation}, par l'image d'\textbf{un seul}
        point.
    \end{itemize}
    Le nombre de points nécessaires mesure la « taille » du groupe ---
    une manière concrète de faire sentir qu'une rotation est « plus
    riche » qu'une translation.
  \end{Rem}
\end{frame}
